\section{Tangent maps, linearizations, and effective transfer}
\label{sec:tangent-effective-v163}
We separate the source of the contact tangent map from the category in
which the boundary is embedded. The elementary linearizations below
also specify the splitting used in the envelope automorphism group.

\begin{lemma}[Framed tangent map and stabilizer]
\label{lem:framed-tangent-v163}
Let $A(x)$ be a regular symmetric pencil, expressed as a symmetric map
from $E$ to $E^*$ with its parameter-line twist trivialized locally.
After splitting a unit block, let $j(x):W[[x]]\to E[[x]]$ be the
inclusion of its orthogonal residual summand, and let $D(x)$ be the
residual diagonal form. On the atlas of ordered pencil coefficients
the contact differential is
\begin{equation}\label{eq:contact-differential-v163}
 (\dot A_0,\dot A_1)\longmapsto
 \left[j(x)^*(\dot A_0+x\dot A_1)j(x)\right]
 \quad\hbox{in }\frac{\Sym^2W^*[[x]]}
 {\{H^{\mathsf t}D+DH:H\in\operatorname{End}(W)[[x]]\}}.
\end{equation}
The restriction to distinct contacts gives the direct sum of these
maps. Local frame changes act by congruence and by the corresponding
line-bundle unit; they identify these modules and differentials.
This formula does not discard the pencil-basis directions.

If a nonsingular member is fixed and $T$ is its self-adjoint pencil
operator, the Lie algebra of its orthogonal stabilizer is
\[
        \{K:KT=TK,\ K^*=-K\}.
\]
For Jordan block sizes $d_{\lambda,i}$ its dimension is
$\sum_{\lambda,i<j}\min(d_{\lambda,i},d_{\lambda,j})$. Hence on
this fixed-member framed slice the contact quotient has dimension
\[
 n+\sum_{\lambda,i<j}\min(d_{\lambda,i},d_{\lambda,j}).
\]
This is not asserted to be the dimension of the unframed pencil
quotient by source congruence and parameter reparametrization.
\end{lemma}
\begin{proof}
For a block matrix with invertible upper block $U$ and off-diagonal
block $B$, the inclusion is $j=(-U^{-1}B,1)^{\mathsf t}$. Differentiating
the Schur complement gives $j^*\dot A j$: the terms differentiating
$j$ cancel because its image is orthogonal to the unit block.
Infinitesimal residual frame changes contribute exactly
$H^{\mathsf t}D+DH$, proving \eqref{eq:contact-differential-v163}.
The polarizing line of the quadratic form has not disappeared. If
$\ell$ is the tautological line at the pencil point, the untrivialized
value is in $\Sym^2W^*\otimes\ell^*$, with the same congruence
quotient; choosing a frame of $\ell$ gives the displayed formula.
Changing it multiplies both numerator and congruence image by the
same unit.

The ordered-basis atlas has tangent space $\Hom(\C^2,\Sym^2E^*)$.
At its image pencil $R$, the sequence forgetting the basis is
\[
 0\longrightarrow\operatorname{End}(R)
 \longrightarrow\Hom(R,\Sym^2E^*)
 \longrightarrow\Hom(R,\Sym^2E^*/R)\longrightarrow0.
\]
The rightmost term is the Grassmannian tangent space. A splitting
specifies a local basis section; a different splitting may change
the parameter-fixed contact coordinates. Thus the differential is
first formulated on the framed atlas. The comparison theorem
retains these choices and descends by its transition maps, instead
of silently treating the four basis directions as zero contact
vectors. Source-frame directions contribute the infinitesimal
congruence image. The block computation in
Theorem~\ref{thm:ambient-versal-v162} proves surjectivity on this
atlas, so it proves formal smoothness of its completed contact map.

For the stabilizer, decompose the self-adjoint operator into
orthogonal primary cyclic blocks. Intertwiners between distinct
eigenvalues vanish. Between blocks of the same eigenvalue their
space has dimension $\min(d_i,d_j)$, as homomorphisms of
$\C[t]$-modules $\C[t]/(t^{d_i})$ and $\C[t]/(t^{d_j})$.
An off-diagonal intertwiner $K_{ij}$ forces
$K_{ji}=-K_{ij}^*$ and supplies one independent copy. On a single
cyclic block a commuting endomorphism is a polynomial in $T$,
hence is self-adjoint. It is skew-adjoint only when it is zero.
This proves the dimension formula without assuming semisimplicity.
Subtracting the orbit dimension
$n(n-1)/2-\dim\operatorname{Stab}$ from the dimension
$n(n+1)/2$ of self-adjoint operators gives the last formula.
\end{proof}

\begin{proposition}[The specified envelope has a split automorphism extension]
\label{prop:linearized-envelope-v163}
Fix $V$, $h\geq1$, and the graded algebra bundle on $\Pj(V)$ with
pieces $B_h(V)_p\otimes\OO(2p)$, including its multiplication and
power diagram. Its automorphism functor, allowing automorphisms of
$\Pj(V)$, is represented by a finite-type affine complex group $H_h$.
There is a split exact sequence
\[
 1\longrightarrow K_h\longrightarrow H_h\longrightarrow
                          \operatorname{PGL}(V)\longrightarrow1.
\]
The splitting is the one induced by the explicit source-normalized
linearization, not a general assertion about automorphism extensions.
The fppf stack of forms of this fixed object is the algebraic
classifying stack $BH_h$. This statement permits twisted projective
source bundles but does not classify all graded algebra bundles
of those ranks, or all raw failure algebras.
\end{proposition}
\begin{proof}
Start with the natural $\operatorname{GL}(V)$ action on the apolar
quotient and on the tautological bundle. A scalar $t$ acts on
$B_h(V)_p$ by $t^{2p}$ and on $\OO(2p)$ by $t^{-2p}$. It therefore
acts trivially on their tensor product and on its multiplication.
The action factors through $\operatorname{PGL}(V)$, giving a group
homomorphism into automorphisms of the specified object whose
projection to source automorphisms is the identity. This is the
required splitting and defines the conjugation action on $K_h$.
Without this chosen linearization there would only be an extension,
and a semidirect-product notation would require justification.

An automorphism over the identity of $\Pj(V)$ is, in each degree,
a constant invertible matrix, since
$H^0(\Pj(V),\mathcal End(B_h(V)_p\otimes\OO(2p)))
=\operatorname{End}(B_h(V)_p)$. There are finitely many nonzero
degrees. Preservation of the unit and multiplication is a finite
collection of polynomial equations in these matrices and their
inverses. The resulting closed subgroup $K_h$ is affine of finite
type, and $H_h=K_h\rtimes\operatorname{PGL}(V)$ has the same
properties. An isomorphism torsor for fppf-local forms of this one
object is an $H_h$-torsor, and effective descent gives the inverse
construction. Thus the category of forms is precisely $BH_h$.
Affine finite-type groups over $\C$ are smooth, and the usual
classifying-stack atlas establishes the asserted algebraicity.
\end{proof}

\begin{corollary}[Complete boundary fibres on the effective failure image]
\label{cor:effective-full-fibres-v163}
On the effectively rigidified sharp pencil-failure image, families
satisfying \eqref{eq:first-boundary-class-v163} have the complete
boundary fibres of Theorem~\ref{thm:product-fibres-v163}, including
their nilpotent structures, component dimensions and intersections.
The exponent $h=1$ defines the envelope without an additional integer
choice; other $h\geq1$ have the same graph. The assertion is a
functorial consequence of the sharp inverse on effective families.
It is not an internal boundary construction on a closed algebra
before reconstruction and is not a statement about the full raw
Artin-algebra deformation stack.
\end{corollary}
\begin{proof}
Use the family inverse to recover the oriented source and pencil,
then the source-normalized envelope to obtain its multiplication
ideals. On the corank-two rank open the relevant graph is determined
by $I_{n-1}$; the universal power identity gives the same graph for
every $h$. The completed comparison and the full-fibre theorem
identify its embedded family. Their transition maps commute with
the effective equivalence and the source linearization, so all the
scheme structures descend, including the local ideal
$(eA,eB,e^2C)$ and its products. This proves the claim in the stated
category and no larger one.
\end{proof}
